\section{Erro interpolatório ótimo em árvores estritamente diádicas}
\label{sec:dyadic-asymptotic}

Fixe primeiro o objeto cuja constante será determinada. Para uma célula
diádica $J=[a,b]$ e $f\in C(J)$, denote por $P_Jf$ a reta que interpola
$f(a)$ e $f(b)$ e ponha
\begin{equation}
 e_f(J)=\|f-P_Jf\|_{\infty,J}.
 \label{eq:local-interpolation-error}
\end{equation}
Se $\Lambda$ é o conjunto de folhas de uma árvore diádica finita com raiz
$K$, as retas $P_Jf$, $J\in\Lambda$, concordam nos nós compartilhados e
definem o interpolante contínuo $P_\Lambda f$. Defina
\begin{equation}
 E_N^{\rm dia}(f)_\infty
 =\min_{\substack{\Lambda\text{ árvore diádica}\\ |\Lambda|\le N}}
   \|f-P_\Lambda f\|_{\infty,K}
 =\min_{\substack{\Lambda\text{ árvore diádica}\\ |\Lambda|\le N}}
   \max_{J\in\Lambda}e_f(J).
 \label{eq:optimal-dyadic-interpolation}
\end{equation}
O mínimo existe: há apenas um número finito de árvores binárias ordenadas e
cheias com no máximo $N$ folhas, e uma árvore com $n$ folhas tem profundidade
no máximo $n-1$.

O sobrescrito ``dia'' distingue \eqref{eq:optimal-dyadic-interpolation} do
erro truncado \eqref{eq:truncated-tree-error}. O primeiro otimiza a
\emph{árvore} e fixa em cada folha a reta interpoladora real; o segundo
otimiza simultaneamente árvore e rótulos racionais de altura no máximo $H$.
As constantes abaixo pertencem ao primeiro problema.
Elas serão comparadas com a referência de nós livres do
\cref{thm:free-knots}; a questão aqui é precisamente o custo adicional de
restringir todos os nós a uma árvore diádica.

\subsection{A árvore de limiar é exatamente mínima}

Para $\varepsilon>0$, construa $T_\varepsilon(f)$ recursivamente: uma célula
$J$ é mantida como folha quando $e_f(J)\le\varepsilon$ e é substituída por
seus dois filhos quando $e_f(J)>\varepsilon$. Denote por
\begin{equation}
 n_f(\varepsilon)=|\Lambda(T_\varepsilon(f))|
 \label{eq:threshold-count}
\end{equation}
o número de folhas. Essa regra usa o erro local \emph{exato}; o refinamento
certificado da \cref{sec:adaptive} usa o majorante $|J|^2M(J)/8$.

\begin{theorem}[minimalidade exata do limiar]
\label{thm:threshold-optimality}
Se $f\in C(K)$, então $T_\varepsilon(f)$ é finita e é a única árvore
menos refinada (\emph{coarsest}) cuja interpolação tem erro no máximo $\varepsilon$. Em particular,
ela possui o menor número de folhas entre todas as árvores viáveis e
\begin{equation}
 E_N^{\rm dia}(f)_\infty\le\varepsilon
 \quad\Longleftrightarrow\quad
 n_f(\varepsilon)\le N,
 \qquad
 E_N^{\rm dia}(f)_\infty
 =\inf\{\varepsilon>0:n_f(\varepsilon)\le N\}.
 \label{eq:threshold-inversion}
\end{equation}
\end{theorem}

\begin{proof}
Se $\omega_f$ é o módulo de continuidade de $f$, a combinação convexa que
define $P_Jf$ dá
\[
 e_f(J)\le\omega_f(|J|).
\]
Logo todas as células de profundidade suficientemente grande satisfazem o
limiar, e a recursão termina. Quando $f\in C^2$, vale ainda a cota mais
precisa
\begin{equation}
 e_f(J)\le \frac{|J|^2}{8}\|f''\|_{\infty,J}.
 \label{eq:local-upper-basic}
\end{equation}

Se $e_f(J)>\varepsilon$, nenhuma árvore viável pode manter $J$ como folha;
ela é obrigada a subdividi-la. Se $e_f(J)\le\varepsilon$, manter $J$ já é
viável e usa menos folhas que qualquer subdivisão própria. Aplicar essa
dicotomia recursivamente prova, por indução na subárvore finita produzida,
que toda árvore viável refina $T_\varepsilon(f)$. Isso dá unicidade da árvore menos
refinada e minimalidade do número de folhas. A equivalência e a inversão seguem da
definição de \eqref{eq:optimal-dyadic-interpolation}.
\end{proof}

O \cref{thm:threshold-optimality} aplica-se ao funcional uniforme local: a
poda por limiar é exata porque a
viabilidade em norma uniforme é a conjunção dos testes locais
$e_f(J)\le\varepsilon$, e toda falha obriga a mesma subdivisão. Esse problema
escalar é uma especialização da recorrência min--max de
\cref{sec:variable-dp} na qual os
rótulos das folhas já foram fixados pelos dados interpolados. A literatura de
aproximação em árvores enfatiza que, para objetivos mais gerais, árvores
ótimas e algoritmos adaptativos quase ótimos são noções distintas
\citep{cohen2001,binev2004}.

\subsection{Lema local e funcional de curvatura}

Escreva
\[
 \omega_{f''}(t)=\sup_{|x-y|\le t}|f''(x)-f''(y)|,
 \qquad
 I(f)=\int_K\sqrt{|f''(x)|}\,\dd x.
\]

\begin{lemma}[expansão local uniforme]
\label{lem:local-asymptotic-error}
Para toda célula $J\subseteq K$ e todo $x\in J$,
\begin{equation}
 \left|e_f(J)-\frac{|f''(x)|}{8}|J|^2\right|
 \le \frac{\omega_{f''}(|J|)}{8}|J|^2.
 \label{eq:local-asymptotic-error}
\end{equation}
\end{lemma}

\begin{proof}
Seja $q_x$ o polinômio quadrático com $q_x''\equiv f''(x)$ obtido do jato de
$f$ em $x$. Como a interpolação reproduz funções afins,
\[
 \|q_x-P_Jq_x\|_{\infty,J}=\frac{|f''(x)|}{8}|J|^2.
\]
Aplicar \eqref{eq:local-upper-basic} a $f-q_x$ dá
\[
 \|(f-P_Jf)-(q_x-P_Jq_x)\|_{\infty,J}
 \le\frac{|J|^2}{8}\|f''-f''(x)\|_{\infty,J}.
\]
A desigualdade triangular inversa para normas conclui.
\end{proof}

\subsection{O envelope assintótico geral}

\begin{theorem}[envelope diádico universal e agudo]
\label{thm:general-dyadic-envelope}
Se $f\in C^2(K)$ e $I(f)>0$, então
\begin{align}
 \frac{I(f)}{\sqrt8}
 &\le \liminf_{\varepsilon\downarrow0}
       \sqrt\varepsilon\,n_f(\varepsilon)
 \le \limsup_{\varepsilon\downarrow0}
       \sqrt\varepsilon\,n_f(\varepsilon)
 \le \frac{I(f)}{\sqrt2},
 \label{eq:threshold-envelope}\\
 \frac{I(f)^2}{8}
 &\le \liminf_{N\to\infty}N^2E_N^{\rm dia}(f)_\infty
 \le \limsup_{N\to\infty}N^2E_N^{\rm dia}(f)_\infty
 \le \frac{I(f)^2}{2}.
 \label{eq:budget-envelope}
\end{align}
As afirmações continuam válidas se $f''$ possui zeros ou muda de sinal. Se
$I(f)=0$, então $f$ é afim e ambos os erros são identicamente zero.
\end{theorem}

\begin{proof}
Fixe uma partição uniforme diádica $\mathcal Q_q$ de $K$ e ponha
\[
 m_Q=\min_Q|f''|,
 \qquad M_Q=\max_Q|f''|,
 \qquad
 S_q=\sum_{Q\in\mathcal Q_q}|Q|\sqrt{m_Q},
 \quad U_q=\sum_{Q\in\mathcal Q_q}|Q|\sqrt{M_Q}.
\]
Por continuidade de $\sqrt{|f''|}$, $S_q\to I(f)$ e $U_q\to I(f)$.

Considere primeiro $Q$ com $m_Q>0$. O sinal de $f''$ é constante em $Q$.
Para $\varepsilon$ suficientemente pequeno, todos os ancestrais diádicos
próprios de $Q$ têm erro maior que $\varepsilon$: cada um contém um trecho
de curvatura não nula e, portanto, $f$ não é afim nele. Assim, as folhas de
$T_\varepsilon(f)$ que cobrem $Q$ estão contidas em $Q$. No ponto médio de
uma dessas folhas $J$, a fórmula do erro de interpolação fornece
\[
 e_f(J)\ge \frac{m_Q}{8}|J|^2.
\]
Como $e_f(J)\le\varepsilon$, são necessárias pelo menos
$|Q|\sqrt{m_Q}/\sqrt{8\varepsilon}$ folhas para cobrir $Q$. Somando os
$Q$ com $m_Q>0$,
\[
 \liminf_{\varepsilon\downarrow0}
 \sqrt\varepsilon\,n_f(\varepsilon)\ge S_q/\sqrt8.
\]
Fazer $q\to\infty$ dá a primeira desigualdade de
\eqref{eq:threshold-envelope}. Esse argumento também explica por que zeros e
mudanças de sinal não precisam ser excluídos: as células com $m_Q=0$ são
omitidas na soma inferior, e sua contribuição reaparece no limite de Riemann.

Para a cota superior, refine cada $Q\in\mathcal Q_q$ uniformemente até o
primeiro nível $r_Q$ para o qual
\[
 \frac{M_Q|Q|^2}{8}\,4^{-r_Q}\le\varepsilon.
\]
Essa é uma árvore viável por \eqref{eq:local-upper-basic}. Quando
$M_Q|Q|^2/8>\varepsilon$,
\[
 2^{r_Q}<2\sqrt{\frac{M_Q|Q|^2}{8\varepsilon}}
 =\frac{|Q|\sqrt{M_Q}}{\sqrt{2\varepsilon}};
\]
quando a desigualdade falha, usa-se uma única folha. Pela minimalidade do
\cref{thm:threshold-optimality},
\[
 n_f(\varepsilon)
 \le 2^q+\frac{U_q}{\sqrt{2\varepsilon}}.
\]
Primeiro faça $\varepsilon\downarrow0$ e depois $q\to\infty$ para concluir
\eqref{eq:threshold-envelope}.

As partições uniformes e \eqref{eq:local-upper-basic} mostram primeiro que
$E_N^{\rm dia}(f)_\infty\to0$. Se
$\varepsilon_N=E_N^{\rm dia}(f)_\infty$, então
$n_f(\varepsilon_N)\le N$ por \eqref{eq:threshold-inversion}; a cota inferior
acabada de provar dá
$\liminf N^2\varepsilon_N\ge I(f)^2/8$. Para a direção oposta, fixe
$\delta>0$, escolha $q$ com $U_q<I(f)+\delta$ e tome
\[
 \varepsilon=\frac{(I(f)+2\delta)^2}{2N^2}.
\]
Para $N$ grande, a estimativa anterior fornece
$n_f(\varepsilon)\le N$. Logo $E_N^{\rm dia}\le\varepsilon$ pela inversão.
Fazer $N\to\infty$ e depois $\delta\downarrow0$ prova a última desigualdade
de \eqref{eq:budget-envelope}.
\end{proof}

O fator quatro entre os extremos de \eqref{eq:budget-envelope} não é folga
do argumento. A próxima seção mostra que uma única função quadrática atinge
os dois extremos e, entre eles, todos os valores intermediários.
